% Options for packages loaded elsewhere
\PassOptionsToPackage{unicode}{hyperref}
\PassOptionsToPackage{hyphens}{url}
\documentclass[
  11pt,
]{article}
\usepackage{xcolor}
\usepackage[margin=2.54cm]{geometry}
\usepackage{amsmath,amssymb}
\setcounter{secnumdepth}{-\maxdimen} % remove section numbering
\usepackage{iftex}
\ifPDFTeX
  \usepackage[T1]{fontenc}
  \usepackage[utf8]{inputenc}
  \usepackage{textcomp} % provide euro and other symbols
\else % if luatex or xetex
  \usepackage{unicode-math} % this also loads fontspec
  \defaultfontfeatures{Scale=MatchLowercase}
  \defaultfontfeatures[\rmfamily]{Ligatures=TeX,Scale=1}
\fi
\usepackage{lmodern}
\ifPDFTeX\else
  % xetex/luatex font selection
\fi
% Use upquote if available, for straight quotes in verbatim environments
\IfFileExists{upquote.sty}{\usepackage{upquote}}{}
\IfFileExists{microtype.sty}{% use microtype if available
  \usepackage[]{microtype}
  \UseMicrotypeSet[protrusion]{basicmath} % disable protrusion for tt fonts
}{}
\makeatletter
\@ifundefined{KOMAClassName}{% if non-KOMA class
  \IfFileExists{parskip.sty}{%
    \usepackage{parskip}
  }{% else
    \setlength{\parindent}{0pt}
    \setlength{\parskip}{6pt plus 2pt minus 1pt}}
}{% if KOMA class
  \KOMAoptions{parskip=half}}
\makeatother
\setlength{\emergencystretch}{3em} % prevent overfull lines
\providecommand{\tightlist}{%
  \setlength{\itemsep}{0pt}\setlength{\parskip}{0pt}}
\usepackage{graphicx}
\graphicspath{{../}}
\usepackage{booktabs}
\usepackage{tabularx}
\usepackage[table]{xcolor}
\usepackage{array}
\usepackage{float}
\usepackage{bookmark}
\IfFileExists{xurl.sty}{\usepackage{xurl}}{} % add URL line breaks if available
\urlstyle{same}
\hypersetup{
  hidelinks,
  pdfcreator={LaTeX via pandoc}}

\author{}
\date{}

\begin{document}

\section{EBU6307 Stereo Depth Mini Project
Report}\label{ebu6307-stereo-depth-mini-project-report}

\textbf{Student name:} Zixi Yu\\
\textbf{Student number:} 231223210\\
\textbf{Submission package:} EBU6307\_ZIXI\_YU\_231223210

\newsavebox{\exampleleftbox}
\newcolumntype{L}{>{\raggedright\arraybackslash}X}
\newcolumntype{C}{>{\centering\arraybackslash}X}
\newcommand{\theadl}[1]{\multicolumn{1}{l}{\textbf{#1}}}
\newcommand{\theadc}[1]{\multicolumn{1}{c}{\textbf{#1}}}
\newcommand{\reportdoubletoprule}{%
\specialrule{1pt}{0pt}{1pt}%
\toprule
}
\newcommand{\exampletriple}[6]{%
\begin{center}%
\sbox{\exampleleftbox}{%
\begin{minipage}[t]{0.58\linewidth}
\centering
\vspace{0pt}
\includegraphics[width=0.98\linewidth]{#1}\\[-0.2em]
{\footnotesize #2}

\vspace{0.6em}
\includegraphics[width=0.98\linewidth]{#3}\\[-0.2em]
{\footnotesize #4}
\end{minipage}%
}%
\usebox{\exampleleftbox}\hspace{0.015\linewidth}
\begin{minipage}[t]{0.36\linewidth}
\centering
\vspace{0pt}
\includegraphics[height=\dimexpr\ht\exampleleftbox+\dp\exampleleftbox-2.2em\relax]{#5}\\[-0.2em]
{\footnotesize #6}
\end{minipage}
\end{center}%
}
\newcommand{\annotatedexampleimage}[4][\linewidth]{%
\begin{minipage}[t]{#1}
\centering
\includegraphics[width=\linewidth]{#2}\\[-0.2em]
{\footnotesize #4}
\end{minipage}
}
\newcommand{\exampletriptychrow}[9]{%
\annotatedexampleimage[0.32\linewidth]{#1}{#2}{#3}\hfill
\annotatedexampleimage[0.32\linewidth]{#4}{#5}{#6}\hfill
\annotatedexampleimage[0.32\linewidth]{#7}{#8}{#9}\par
}
\newcommand{\appendixpipelinefigure}[1]{%
\begin{center}
\includegraphics[angle=90,totalheight=0.9\textheight]{#1}
\end{center}
}
\newcommand{\siftexamplerow}[9]{%
\begin{minipage}[t]{0.98\linewidth}
\centering
\annotatedexampleimage[0.49\linewidth]{#1}{#2}{#3}\hfill
\annotatedexampleimage[0.49\linewidth]{#4}{#5}{#6}\\[0.45em]
\annotatedexampleimage[\linewidth]{#7}{#8}{#9}
\end{minipage}\par
}
\newcommand{\siftexamplerowsmall}[9]{%
\begin{minipage}[t]{0.72\linewidth}
\centering
\annotatedexampleimage[0.49\linewidth]{#1}{#2}{#3}\hfill
\annotatedexampleimage[0.49\linewidth]{#4}{#5}{#6}\\[0.45em]
\annotatedexampleimage[\linewidth]{#7}{#8}{#9}
\end{minipage}\par
}

\subsection{1. Project Overview}\label{project-overview}

This report documents the four objectives required by the EBU6307 Mini
Project. All evidence in this report is taken directly from the
submitted package \texttt{EBU6307\_ZIXI\_YU\_231223210}, where the
implementations are stored in \texttt{codes/} and the deliverables are
stored in \texttt{results/}. The experiments cover 24 official
Middlebury scenes and follow the PDF-required structure for each
objective: pipeline figure, three examples, quantitative evaluation, and
a short discussion.

The project progresses from classical geometric synthesis to local
feature engineering, then to dense stereo matching, and finally to a
lightweight Vision Transformer style stereo model. This organization
makes it possible to compare how much accuracy and robustness are gained
by moving from hand-crafted features to learned token descriptors.

\begin{table}[H]
\centering
\caption{\textbf{Project objectives and deliverables.} Overview of the four mini-project objectives, their main entry scripts, required result folders, and the core quantitative metric used for evaluation.}
\label{tab:project_overview_en}
\vspace{1mm}
\small
\renewcommand{\arraystretch}{1.03}
\setlength{\tabcolsep}{5pt}
\begin{tabularx}{\columnwidth}{L C L C}
\reportdoubletoprule
\theadl{Objective} & \theadc{Main script} & \theadl{Required result folders} & \theadc{Core metric} \\
\midrule
O1 Synthetic stereo generation & \texttt{codes/o1.py} & results/O1a\_*, results/O1b\_*, results/O1c\_* & SSIM \\
O2 Manual SIFT repeatability & \texttt{codes/o2.py} & results/O2a\_*, results/O2b\_*, results/O2c\_* & Repeatability \\
O3 SIFT-guided stereo depth & \texttt{codes/o3.py} & results/O3a\_*, results/O3b\_*, results/O3c\_* & MAE / RMSE / Bad-1px \\
O4 Transformer stereo depth & \texttt{codes/o4.py} & results/O4a\_*, results/O4b\_*, results/O4c\_* & 5-fold CV MAE / RMSE / Bad-1px \\
\bottomrule
\specialrule{1pt}{1pt}{0pt}
\end{tabularx}
\vspace{1mm}
\end{table}

\subsection{2. Objective 1: Image Reprojection for Synthetic Stereo
Generation}\label{objective-1-image-reprojection-for-synthetic-stereo-generation}

\subsubsection{2.1 a. Pipeline Figure}\label{a.-pipeline-figure}

\begin{center}
\includegraphics[width=\linewidth]{EBU6307_ZIXI_YU_231223210/results/O1a_synthetic_data/syn_pipeline.jpg}
\end{center}

Objective 1 generates a synthetic right view from the original left
image and the provided left-view disparity map. The implementation first
cleans the reference disparity, applies only light smoothing, and then
forward-projects the left RGB image to a synthetic target view. Because
closer pixels are projected after farther pixels, the projection
approximates stereo occlusion ordering instead of simply shifting the
whole image by a fixed amount. Holes caused by occlusion and discrete
disparity are repaired in three stages: 4-neighbour diffusion, linear
interpolation over unresolved regions, and a short final diffusion pass
to blend the boundaries.

This stage is implemented in \texttt{codes/o1.py} and writes
dataset-style scene folders so later objectives can inspect or reuse the
generated stereo pairs. Each scene folder contains \texttt{im0.png},
synthetic \texttt{im1.png}, \texttt{disp0.pfm}, and copied calibration
metadata where available. The quantitative score is the SSIM between the
true right image and the synthesized right image.

\subsubsection{2.2 b. Three Generated
Examples}\label{b.-three-generated-examples}

\exampletriple{EBU6307_ZIXI_YU_231223210/results/O1b_synthetic_data/artroom2/im1.png}{(1) artroom2 synthetic right view}{EBU6307_ZIXI_YU_231223210/results/O1b_synthetic_data/chess3/im1.png}{(2) chess3 synthetic right view}{EBU6307_ZIXI_YU_231223210/results/O1b_synthetic_data/ladder1/im1.png}{(3) ladder1 synthetic right view}

\textbf{Example 1:} \texttt{artroom2} synthetic right view, SSIM =
\textbf{0.8427}. The large foreground flower and cabinet structures
remain visually coherent after reprojection.\\
\textbf{Example 2:} \texttt{chess3} synthetic right view, SSIM =
\textbf{0.7602}. Repeated edges and thin structures make the
reprojection harder, so the score is lower than \texttt{artroom2}.\\
\textbf{Example 3:} \texttt{ladder1} synthetic right view, SSIM =
\textbf{0.8005}. The result preserves the main geometry well, especially
on the ladder rails and background planes.

\subsubsection{2.3 c.~Quantitative
Evaluation}\label{c.-quantitative-evaluation}

All 24 scenes were synthesized with a horizontal shift setting of 8
pixels, and the final evaluation was recorded in
\texttt{results/O1c\_synthetic\_data/SSIM.csv}.

\begin{table}[H]
\centering
\caption{\textbf{Objective 1 summary statistics.} Scene-level SSIM results over all 24 evaluated Middlebury scenes for synthetic stereo generation.}
\label{tab:o1_stats_en}
\vspace{1mm}
\small
\renewcommand{\arraystretch}{1.03}
\setlength{\tabcolsep}{5pt}
\begin{tabular}{c c}
\reportdoubletoprule
\theadc{Statistic} & \theadc{Value} \\
\midrule
Number of scenes & 24 \\
Mean SSIM & 0.7796 \\
Median SSIM & 0.7758 \\
Best scene & \texttt{bandsaw2} with 0.9159 \\
Worst scene & \texttt{skates1} with 0.6562 \\
\bottomrule
\specialrule{1pt}{1pt}{0pt}
\end{tabular}
\vspace{1mm}
\end{table}

\begin{table}[H]
\centering
\caption{\textbf{Objective 1 selected examples.} SSIM values corresponding to the three visual examples shown for synthetic stereo generation.}
\label{tab:o1_examples_en}
\vspace{1mm}
\small
\renewcommand{\arraystretch}{1.03}
\setlength{\tabcolsep}{5pt}
\begin{tabular}{c c}
\reportdoubletoprule
\theadc{Selected example scene} & \theadc{SSIM} \\
\midrule
artroom2 & 0.8427 \\
chess3 & 0.7602 \\
ladder1 & 0.8005 \\
\bottomrule
\specialrule{1pt}{1pt}{0pt}
\end{tabular}
\vspace{1mm}
\end{table}

\subsubsection{2.4 d.~Discussion}\label{d.-discussion}

\begin{itemize}
\tightlist
\item
  The code for O1 directly reflects the pipeline figure: disparity
  cleaning, forward projection, hole filling, and SSIM evaluation are
  all implemented explicitly in \texttt{codes/o1.py}.
\item
  The synthetic stereo generation works reliably on all 24 scenes and
  produces the PDF-required outputs in the correct folder structure.
\item
  The mean SSIM of 0.7796 shows that the generated right views preserve
  most of the scene structure, even though they are produced without
  learning.
\item
  \texttt{artroom2} performs well because the large surfaces and depth
  layers are easy to reproject consistently.
\item
  \texttt{chess3} is more difficult because repeated edges and fine
  structures amplify small projection and interpolation errors.
\item
  The weakest scenes, such as \texttt{skates1}, likely suffer from
  stronger occlusion boundaries and fine foreground detail, which make
  forward warping less stable.
\item
  A major strength of this implementation is that it keeps the true
  disparity geometry rather than compressing the disparity range into a
  toy visualization-only shift.
\item
  The main limitation is that forward projection still creates visible
  missing regions around sharp depth transitions, and these regions must
  be repaired heuristically.
\item
  A clear next step would be sub-pixel splatting or backward warping
  with explicit visibility reasoning to reduce tearing and improve
  structural similarity.
\end{itemize}

\subsection{3. Objective 2: SIFT Feature Descriptor Generation and
Repeatability
Evaluation}\label{objective-2-sift-feature-descriptor-generation-and-repeatability-evaluation}

\subsubsection{3.1 a. Pipeline Figure}\label{a.-pipeline-figure-1}

\begin{center}
\includegraphics[width=\linewidth]{EBU6307_ZIXI_YU_231223210/results/O2a_sift/sift_pipeline.jpg}
\end{center}

An enlarged rotated copy is also provided in the appendix for easier
reading.

Objective 2 evaluates a hand-written SIFT pipeline under deterministic
random image transformations. For each scene, the original image is
transformed by a known affine, rotation, or intensity perturbation. The
detector in \texttt{codes/o2.py} does not call OpenCV SIFT; instead, it
constructs a Gaussian scale space, forms the DoG pyramid, searches for
26-neighbour extrema, refines their location using 3D Taylor expansion,
rejects low-contrast and edge-like responses, assigns dominant
orientations, and builds standard 128-dimensional descriptors from
rotated 4 x 4 x 8 histograms.

Descriptor matching is performed by Euclidean KNN followed by Lowe ratio
filtering. A match is counted as repeatable only if the matched keypoint
agrees with the known transformation homography within a 4-pixel
threshold. This makes the evaluation stricter than simple descriptor
similarity and better aligned with the real notion of geometric
repeatability.

\subsubsection{3.2 b. Three Generated
Examples}\label{b.-three-generated-examples-1}

\begin{center}
\siftexamplerow{EBU6307_ZIXI_YU_231223210/results/O2b_sift/artroom2/im0_keypoints.png}{1a}{(1a) im0 keypoints}{EBU6307_ZIXI_YU_231223210/results/O2b_sift/artroom2/im1_keypoints.png}{1b}{(1b) im1 keypoints}{EBU6307_ZIXI_YU_231223210/results/O2b_sift/artroom2/sift_matches.png}{1c}{(1c) matches}
\vspace{0.6em}

\siftexamplerow{EBU6307_ZIXI_YU_231223210/results/O2b_sift/chess3/im0_keypoints.png}{2a}{(2a) im0 keypoints}{EBU6307_ZIXI_YU_231223210/results/O2b_sift/chess3/im1_keypoints.png}{2b}{(2b) im1 keypoints}{EBU6307_ZIXI_YU_231223210/results/O2b_sift/chess3/sift_matches.png}{2c}{(2c) matches}
\vspace{0.6em}

\siftexamplerowsmall{EBU6307_ZIXI_YU_231223210/results/O2b_sift/ladder1/im0_keypoints.png}{3a}{(3a) im0 keypoints}{EBU6307_ZIXI_YU_231223210/results/O2b_sift/ladder1/im1_keypoints.png}{3b}{(3b) im1 keypoints}{EBU6307_ZIXI_YU_231223210/results/O2b_sift/ladder1/sift_matches.png}{3c}{(3c) matches}
\end{center}

\textbf{Example 1:} \texttt{artroom2}, transform family =
\textbf{affine}, repeatability = \textbf{0.7413}. The im0/im1 keypoint
overlays are dense, and the match visualization remains mostly
geometrically consistent.\\
\textbf{Example 2:} \texttt{chess3}, transform family = \textbf{affine},
repeatability = \textbf{0.7013}. The detector remains reasonably stable
in both keypoint maps despite repeated indoor structure.\\
\textbf{Example 3:} \texttt{ladder1}, transform family =
\textbf{rotation}, repeatability = \textbf{0.6400}. Thin elongated
structures are still visible in the keypoint overlays, but rotational
variation is more challenging than simple intensity changes.

\subsubsection{3.3 c.~Quantitative
Evaluation}\label{c.-quantitative-evaluation-1}

The complete scene-level metrics are stored in
\texttt{results/O2c\_sift/metrics.csv}.

\begin{table}[H]
\centering
\caption{\textbf{Objective 2 summary statistics.} Scene-level repeatability results over all 24 transformed scenes.}
\label{tab:o2_stats_en}
\vspace{1mm}
\small
\renewcommand{\arraystretch}{1.03}
\setlength{\tabcolsep}{5pt}
\begin{tabular}{c c}
\reportdoubletoprule
\theadc{Statistic} & \theadc{Value} \\
\midrule
Number of scenes & 24 \\
Mean repeatability & 0.6506 \\
Median repeatability & 0.6319 \\
Best scene & \texttt{artroom1} with 0.8313 \\
Worst scene & \texttt{skates1} with 0.4662 \\
\bottomrule
\specialrule{1pt}{1pt}{0pt}
\end{tabular}
\vspace{1mm}
\end{table}

\begin{table}[H]
\centering
\caption{\textbf{Objective 2 repeatability by transform family.} Mean repeatability grouped by intensity, affine, and rotation perturbations.}
\label{tab:o2_family_en}
\vspace{1mm}
\small
\renewcommand{\arraystretch}{1.03}
\setlength{\tabcolsep}{5pt}
\begin{tabular}{c c}
\reportdoubletoprule
\theadc{Transform family} & \theadc{Mean repeatability} \\
\midrule
intensity & 0.7225 \\
affine & 0.6429 \\
rotation & 0.5864 \\
\bottomrule
\specialrule{1pt}{1pt}{0pt}
\end{tabular}
\vspace{1mm}
\end{table}

\begin{table}[H]
\centering
\caption{\textbf{Objective 2 selected examples.} Repeatability values for the three visual examples shown in the SIFT evaluation section.}
\label{tab:o2_examples_en}
\vspace{1mm}
\small
\renewcommand{\arraystretch}{1.03}
\setlength{\tabcolsep}{5pt}
\begin{tabular}{c c c}
\reportdoubletoprule
\theadc{Selected example scene} & \theadc{Transform family} & \theadc{Repeatability} \\
\midrule
artroom2 & affine & 0.7413 \\
chess3 & affine & 0.7013 \\
ladder1 & rotation & 0.6400 \\
\bottomrule
\specialrule{1pt}{1pt}{0pt}
\end{tabular}
\vspace{1mm}
\end{table}

\subsubsection{3.4 d.~Discussion}\label{d.-discussion-1}

\begin{itemize}
\tightlist
\item
  O2 demonstrates that the manual implementation in \texttt{codes/o2.py}
  successfully reproduces the main stages of the classical SIFT
  algorithm rather than delegating the task to a black-box library.
\item
  The required deliverables were produced correctly: pipeline figure,
  three example visualizations, and the repeatability CSV table.
\item
  The overall mean repeatability of 0.6506 is solid for a fully
  hand-crafted detector and descriptor evaluated under multiple
  transformation families.
\item
  Intensity changes are handled best, with a mean repeatability of
  0.7225, which indicates that the normalization and orientation
  assignment stages are working as intended.
\item
  Rotational perturbations are the hardest case, with a mean
  repeatability of 0.5864, showing that viewpoint and alignment changes
  still reduce descriptor stability.
\item
  \texttt{artroom2} is a good example because the cluttered scene
  contains many stable corners and textured regions that produce
  consistent matches.
\item
  \texttt{skates1} is the weakest scene, suggesting that reflective or
  low-repeatability structures still challenge the hand-crafted
  descriptor.
\item
  From a code reflection perspective, the detector is explainable and
  educational, but it is also more sensitive to parameter choices than a
  heavily optimized industrial implementation.
\item
  The next improvement would be sub-bin spatial interpolation and
  stronger descriptor pooling around difficult rotations, which should
  increase stability without abandoning the manual implementation.
\end{itemize}

\subsection{4. Objective 3: Stereo Depth Estimation Based on
SIFT}\label{objective-3-stereo-depth-estimation-based-on-sift}

\subsubsection{4.1 a. Pipeline Figure}\label{a.-pipeline-figure-2}

\begin{center}
\includegraphics[width=\linewidth]{EBU6307_ZIXI_YU_231223210/results/O3a_disparity/dep_pipeline.jpg}
\end{center}

An enlarged rotated copy is also provided in the appendix for easier
reading.

Objective 3 turns sparse SIFT evidence into dense stereo depth
estimation. The implementation first detects manual SIFT features on the
left and right views, performs descriptor KNN matching, applies ratio,
mutual, and stereo-geometry filtering, and converts the reliable
correspondences into a sparse disparity prior. This prior is then
injected into a dense census-plus-gradient matching pipeline with
scene-adaptive disparity bounds.

The dense stage uses eight-direction Semi-Global Matching to aggregate
the cost volume. The result is refined by left-right consistency, local
support masking, speckle removal, median filtering, joint weighted
median filtering, and small gap filling. The final depth output is saved
as \texttt{disp0.pfm}, while \texttt{disp0.png} provides a colorized
preview for the report.

\subsubsection{4.2 b. Three Generated
Examples}\label{b.-three-generated-examples-2}

\exampletriple{EBU6307_ZIXI_YU_231223210/results/O3b_disparity/artroom2/disp0.png}{(1) artroom2 disparity}{EBU6307_ZIXI_YU_231223210/results/O3b_disparity/chess3/disp0.png}{(2) chess3 disparity}{EBU6307_ZIXI_YU_231223210/results/O3b_disparity/ladder1/disp0.png}{(3) ladder1 disparity}

\textbf{Example 1:} \texttt{artroom2}, MAE = \textbf{2.0799}, RMSE =
\textbf{5.7954}, Bad-1px = \textbf{0.4266}. The flower and cabinet
geometry are recovered with visible object boundaries, although some
invalid regions remain black.\\
\textbf{Example 2:} \texttt{chess3}, MAE = \textbf{1.7968}, RMSE =
\textbf{6.4951}, Bad-1px = \textbf{0.2930}. The method handles the scene
relatively well despite challenging structure.\\
\textbf{Example 3:} \texttt{ladder1}, MAE = \textbf{2.1101}, RMSE =
\textbf{3.8366}, Bad-1px = \textbf{0.4802}. The thin ladder rails are
preserved, but the strict consistency filtering leaves holes around
narrow structures.

\subsubsection{4.3 c.~Quantitative
Evaluation}\label{c.-quantitative-evaluation-2}

The full scene-level results are stored in
\texttt{results/O3c\_disparity/metrics.csv}.

\begin{table}[H]
\centering
\caption{\textbf{Objective 3 summary statistics.} Dense disparity evaluation results over all 24 benchmark scenes.}
\label{tab:o3_stats_en}
\vspace{1mm}
\small
\renewcommand{\arraystretch}{1.03}
\setlength{\tabcolsep}{5pt}
\begin{tabular}{c c}
\reportdoubletoprule
\theadc{Statistic} & \theadc{Value} \\
\midrule
Number of scenes & 24 \\
Mean MAE & 2.3357 \\
Mean RMSE & 6.8212 \\
Mean Bad-1px & 0.3881 \\
Median MAE & 2.0749 \\
Best scene & \texttt{skates2} with MAE 1.2164 \\
Worst scene & \texttt{podium1} with MAE 7.3048 \\
\bottomrule
\specialrule{1pt}{1pt}{0pt}
\end{tabular}
\vspace{1mm}
\end{table}

\begin{table}[H]
\centering
\caption{\textbf{Objective 3 selected examples.} Error metrics corresponding to the three disparity examples shown in the O3 section.}
\label{tab:o3_examples_en}
\vspace{1mm}
\small
\renewcommand{\arraystretch}{1.03}
\setlength{\tabcolsep}{5pt}
\begin{tabular}{c c c c}
\reportdoubletoprule
\theadc{Selected example scene} & \theadc{MAE} & \theadc{RMSE} & \theadc{Bad-1px} \\
\midrule
artroom2 & 2.0799 & 5.7954 & 0.4266 \\
chess3 & 1.7968 & 6.4951 & 0.2930 \\
ladder1 & 2.1101 & 3.8366 & 0.4802 \\
\bottomrule
\specialrule{1pt}{1pt}{0pt}
\end{tabular}
\vspace{1mm}
\end{table}

\subsubsection{4.4 d.~Discussion}\label{d.-discussion-2}

\begin{itemize}
\tightlist
\item
  O3 successfully connects the feature engineering from O2 with a dense
  stereo backend, and the implementation in \texttt{codes/o3.py} clearly
  reflects that hybrid design.
\item
  The pipeline produced the required architecture figure, three
  disparity examples, and a full disparity error table for all scenes.
\item
  With a mean MAE of 2.3357, O3 provides a strong classical baseline for
  the later learned model in O4.
\item
  A key strength is that sparse SIFT matches inject reliable structure
  into the dense matcher, especially in scenes with clear local texture.
\item
  Visually, O3 often preserves sharper fine detail than O4, as seen
  around the flower boundaries in \texttt{artroom2} and the ladder rails
  in \texttt{ladder1}.
\item
  The tradeoff is density: left-right consistency and support filtering
  remove many uncertain pixels, so the preview contains more black
  invalid regions.
\item
  \texttt{podium1} is the clearest failure case, with much higher error
  than the dataset median, showing that some scenes still break the
  matching assumptions.
\item
  From a testing perspective, the scene-level metric table confirms that
  the method ran end-to-end across the whole evaluation set and produced
  formal PFM outputs.
\item
  The next improvement would be stronger occlusion reasoning and
  adaptive post-processing so that fine details can be preserved without
  leaving so many invalid holes.
\end{itemize}

\subsection{5. Objective 4: Stereo Depth Estimation with a Vision
Transformer Style
Model}\label{objective-4-stereo-depth-estimation-with-a-vision-transformer-style-model}

\subsubsection{5.1 a. Pipeline / Architecture
Figure}\label{a.-pipeline-architecture-figure}

\begin{center}
\includegraphics[width=\linewidth]{EBU6307_ZIXI_YU_231223210/results/O4a_transformer/transformer_pipeline.jpg}
\end{center}

An enlarged rotated copy is also provided in the appendix for easier
reading.

Objective 4 replaces hand-crafted SIFT descriptors with trainable token
descriptors and a lightweight Transformer encoder. The implemented model
is a custom \texttt{StereoTokenTransformer} defined in
\texttt{codes/o4.py} and \texttt{codes/o4\_torch.py}, not an
off-the-shelf large ViT backbone. The pipeline converts grayscale stereo
images into 2 x 2 patch tokens, augments them with gradient and local
context channels, maps them into token embeddings, and then uses cosine
similarity to build a token-level cost volume.

The learned token cost volume is regularized by the same general stereo
reasoning as O3, but on the token grid: scene-adaptive disparity
bounds, eight-direction SGM, token left-right consistency with a
confidence margin, anchor-guided token cleanup, weighted-median
correction, and 2 x 2 upsampling back to full resolution. The formal
outputs are produced with 5-fold cross-validation and saved to
\texttt{results/O4b\_transformer/} and
\texttt{results/O4c\_transformer/}.

\subsubsection{5.2 b. Three Generated
Examples}\label{b.-three-generated-examples-3}

\begin{center}
\exampletriptychrow{EBU6307_ZIXI_YU_231223210/results/O4b_transformer/artroom2/disp0_transformer_raw.png}{1a}{(1a) raw}{EBU6307_ZIXI_YU_231223210/results/O4b_transformer/artroom2/disp0_transformer_raw_filtered.png}{1b}{(1b) raw filtered}{EBU6307_ZIXI_YU_231223210/results/O4b_transformer/artroom2/disp0.png}{1c}{(1c) final disparity}
\vspace{0.6em}

\exampletriptychrow{EBU6307_ZIXI_YU_231223210/results/O4b_transformer/chess3/disp0_transformer_raw.png}{2a}{(2a) raw}{EBU6307_ZIXI_YU_231223210/results/O4b_transformer/chess3/disp0_transformer_raw_filtered.png}{2b}{(2b) raw filtered}{EBU6307_ZIXI_YU_231223210/results/O4b_transformer/chess3/disp0.png}{2c}{(2c) final disparity}
\vspace{0.6em}

\exampletriptychrow{EBU6307_ZIXI_YU_231223210/results/O4b_transformer/ladder1/disp0_transformer_raw.png}{3a}{(3a) raw}{EBU6307_ZIXI_YU_231223210/results/O4b_transformer/ladder1/disp0_transformer_raw_filtered.png}{3b}{(3b) raw filtered}{EBU6307_ZIXI_YU_231223210/results/O4b_transformer/ladder1/disp0.png}{3c}{(3c) final disparity}
\end{center}

\textbf{Example 1:} \texttt{artroom2}, test fold = \textbf{1}, MAE =
\textbf{2.2037}, RMSE = \textbf{6.5174}, Bad-1px = \textbf{0.3794}. The
raw token disparity is noisy, the filtered raw map is cleaner, and the
final map is denser and smoother than O3.\\
\textbf{Example 2:} \texttt{chess3}, test fold = \textbf{1}, MAE =
\textbf{5.2626}, RMSE = \textbf{13.1028}, Bad-1px = \textbf{0.3145}.
Repeated structure is already visible as ambiguity in the raw maps and
remains difficult for the learned token matcher after refinement, leading
to a much larger regression error.\\
\textbf{Example 3:} \texttt{ladder1}, test fold = \textbf{0}, MAE =
\textbf{1.9797}, RMSE = \textbf{3.7391}, Bad-1px = \textbf{0.4337}. The
raw and filtered maps already preserve the ladder structure, and the
final output remains denser than O3 while staying close in error.

\subsubsection{5.3 c.~Quantitative
Evaluation}\label{c.-quantitative-evaluation-3}

The complete results are stored in
\texttt{results/O4c\_transformer/metrics.csv}, and the fold-wise
cross-validation summary is stored in
\texttt{results/O4c\_transformer/fold\_summary.csv}.

\begin{table}[H]
\centering
\caption{\textbf{Objective 4 summary statistics.} Overall five-fold cross-validation disparity evaluation results across the 24 benchmark scenes.}
\label{tab:o4_stats_en}
\vspace{1mm}
\small
\renewcommand{\arraystretch}{1.03}
\setlength{\tabcolsep}{5pt}
\begin{tabular}{c c}
\reportdoubletoprule
\theadc{Statistic} & \theadc{Value} \\
\midrule
Number of scenes & 24 \\
Cross-validation folds & 5 \\
Mean MAE & 3.3176 \\
Mean RMSE & 8.1083 \\
Mean Bad-1px & 0.3721 \\
Median MAE & 2.4049 \\
Best scene & \texttt{skates2} with MAE 1.0481 \\
Worst scene & \texttt{podium1} with MAE 12.7761 \\
\bottomrule
\specialrule{1pt}{1pt}{0pt}
\end{tabular}
\vspace{1mm}
\end{table}

\begin{table}[H]
\centering
\caption{\textbf{Objective 4 fold-wise summary.} Mean performance for each cross-validation fold in the Transformer-based stereo evaluation.}
\label{tab:o4_folds_en}
\vspace{1mm}
\small
\renewcommand{\arraystretch}{1.03}
\setlength{\tabcolsep}{5pt}
\begin{tabular}{c c c c c}
\reportdoubletoprule
\theadc{Fold} & \theadc{Scene count} & \theadc{Mean MAE} & \theadc{Mean RMSE} & \theadc{Mean Bad-1px} \\
\midrule
0 & 5 & 3.5480 & 8.9336 & 0.4122 \\
1 & 5 & 5.0527 & 11.0331 & 0.4554 \\
2 & 5 & 2.4753 & 6.7348 & 0.3182 \\
3 & 5 & 2.4646 & 6.1603 & 0.3522 \\
4 & 4 & 2.9799 & 7.5726 & 0.3102 \\
\bottomrule
\specialrule{1pt}{1pt}{0pt}
\end{tabular}
\vspace{1mm}
\end{table}

\begin{table}[H]
\centering
\caption{\textbf{Objective 4 selected examples.} Fold assignment and error metrics for the three visual examples shown in the O4 section.}
\label{tab:o4_examples_en}
\vspace{1mm}
\small
\renewcommand{\arraystretch}{1.03}
\setlength{\tabcolsep}{5pt}
\begin{tabular}{c c c c c}
\reportdoubletoprule
\theadc{Selected example scene} & \theadc{Fold} & \theadc{MAE} & \theadc{RMSE} & \theadc{Bad-1px} \\
\midrule
artroom2 & 1 & 2.2037 & 6.5174 & 0.3794 \\
chess3 & 1 & 5.2626 & 13.1028 & 0.3145 \\
ladder1 & 0 & 1.9797 & 3.7391 & 0.4337 \\
\bottomrule
\specialrule{1pt}{1pt}{0pt}
\end{tabular}
\vspace{1mm}
\end{table}

\subsubsection{5.4 d.~Discussion}\label{d.-discussion-3}

\begin{itemize}
\tightlist
\item
  O4 satisfies the PDF requirement for a Vision Transformer based
  objective by replacing SIFT descriptors with a trainable Transformer
  token encoder and reporting 5-fold cross-validation results.
\item
  The implementation is more ambitious than O3 because it combines
  learned descriptors with structured stereo regularization instead of
  relying on either one alone.
\item
  The strongest qualitative advantage of O4 is output density: compared
  with O3, the disparity maps are usually smoother and contain fewer
  black invalid holes.
\item
  Quantitatively, O4 achieves a slightly better mean Bad-1px than O3
  (0.3721 vs 0.3881), which suggests that many pixels are closer to the
  correct disparity within a 1-pixel tolerance.
\item
  However, O4 is worse than O3 on mean MAE and mean RMSE, showing that
  the learned model still makes larger absolute errors on difficult
  scenes.
\item
  \texttt{chess1}, \texttt{chess2}, \texttt{chess3}, and especially
  \texttt{podium1} show that the current token model is still vulnerable
  to repeated structures and hard scene-specific ambiguity.
\item
  The fold summary also shows visible variance across splits, with fold
  1 clearly harder than the others, so the current model is not yet
  fully stable.
\item
  From a code reflection perspective, the model is compact and
  reproducible, with 478,309 trainable parameters and a clear
  checkpointed 5-fold workflow.
\item
  The next improvement would be to strengthen the token descriptor with
  richer context and stronger confidence-aware filtering so the model
  keeps O4's density advantage while reducing the large-scene outliers.
\end{itemize}

\subsection{6. O3, O4, and DINOv2 Comparison}\label{o3-and-o4-comparison}

Because O4 is designed as the learned alternative to O3, the most
important comparison is between the classical and Transformer-based
stereo systems on the same 24 scenes, and an additional pretrained
DINOv2 test is included to show the effect of stronger pretrained
features.

\begin{table}[H]
\centering
\caption{\textbf{Comparison between O3, O4, and pretrained DINOv2.} Aggregate metric comparison on the shared 24-scene benchmark, with better values highlighted in bold.}
\label{tab:o3_o4_compare_en}
\vspace{1mm}
\small
\renewcommand{\arraystretch}{1.03}
\setlength{\tabcolsep}{4pt}
\begin{tabular}{c c c c c}
\reportdoubletoprule
\theadc{Metric} & \theadc{O3} & \theadc{O4} & \theadc{DINOv2} & \theadc{Best method} \\
\midrule
Mean MAE & 2.3357 & 3.3176 & \textbf{2.1023} & DINOv2 \\
Mean RMSE & 6.8212 & 8.1083 & \textbf{5.5086} & DINOv2 \\
Mean Bad-1px & 0.3881 & 0.3721 & \textbf{0.2729} & DINOv2 \\
Mean valid ratio & 0.5316 & \textbf{0.8612} & 0.2035 & O4 \\
Scene wins by MAE & 3 & 0 & \textbf{21} & DINOv2 \\
Scene wins by RMSE & 4 & 1 & \textbf{19} & DINOv2 \\
Scene wins by Bad-1px & 2 & 0 & \textbf{22} & DINOv2 \\
\bottomrule
\specialrule{1pt}{1pt}{0pt}
\end{tabular}
\vspace{1mm}
\end{table}

This extended comparison shows a more nuanced tradeoff. If only the
reported error on predicted pixels is considered, the additional
pretrained DINOv2 test is clearly the strongest result: it delivers the
lowest mean MAE, RMSE, and Bad-1px, and it wins most scenes under all
three metrics. This suggests that strong pretrained descriptors are very
helpful in the current low-data regime. However, its mean valid ratio is
only 0.2035, far below O4's 0.8612 and O3's 0.5316, so those gains are
achieved on a much smaller subset of pixels. O4 therefore remains the
densest learned result, while O3 remains the stronger higher-coverage
baseline when absolute disparity precision is prioritized. At the same
time, O3 still shows visible holes in low-confidence regions, so its
numerical advantage should also be read with some caution: if those
missing areas were penalized as prediction failures, the effective error
would increase further. Overall, the DINOv2 test confirms the promise
of pretrained learned descriptors, but the current compact Transformer
still needs more training data and stronger local detail recovery, for
example CNN or U-Net style refinement, to combine that accuracy with
dense predictions.

\subsection{7. Conclusion}\label{conclusion}

This mini project successfully delivered all four required objectives in
the requested folder structure. O1 generated synthetic stereo pairs with
good structural similarity, O2 implemented a fully manual SIFT pipeline
with meaningful repeatability scores, O3 produced a strong classical
stereo baseline, and O4 extended the project into a learned
Transformer-based stereo system with 5-fold cross-validation.

The most important overall finding is that the project does not simply
progress from ``traditional'' to ``deep'' methods in a superficial way.
Instead, each stage builds technically on the previous one: O1 provides
controlled image synthesis, O2 validates local feature robustness, O3
converts feature evidence into dense geometry, and O4 tests whether
trainable token descriptors can outperform hand-crafted features. At the
current stage, the additional pretrained DINOv2 test achieves the best
reported MAE, RMSE, and Bad-1px, which shows that strong pretraining is
already valuable under the current small-data setting. However, that
accuracy comes with a very low valid ratio, so it is not yet the best
overall solution when dense prediction quality is also required. O3
still provides a stronger higher-coverage handcrafted baseline, O4
provides the densest learned output, and O3 still leaves holes in some
difficult regions, which means its reported error remains somewhat
optimistic if those missing pixels are not explicitly counted. This
gives a clear roadmap for future work: preserve O4's density advantage,
reuse stronger pretraining such as DINOv2, enlarge the training set
beyond the current expanded data, and add stronger detail-recovery
modules such as CNN or U-Net style refinement so that boundaries and
thin structures can be resolved more accurately while improving
robustness on repeated-pattern and hard-ambiguity scenes.

\clearpage
\appendix
\section{Appendix: O2--O4 Pipeline Figures}

\subsection*{Objective 2 Pipeline Figure}
\appendixpipelinefigure{EBU6307_ZIXI_YU_231223210/results/O2a_sift/sift_pipeline.jpg}

\clearpage
\subsection*{Objective 3 Pipeline Figure}
\appendixpipelinefigure{EBU6307_ZIXI_YU_231223210/results/O3a_disparity/dep_pipeline.jpg}

\clearpage
\subsection*{Objective 4 Pipeline / Architecture Figure}
\appendixpipelinefigure{EBU6307_ZIXI_YU_231223210/results/O4a_transformer/transformer_pipeline.jpg}

\end{document}
